\chapter*{Agradecimientos}

\noindent Comencé la carrera formalmente en el año 2022, pero en verdad fue mucho antes. Todo empezó cuando en 2017 comencé a participar en la Olimpíada Matemática Argentina, y especialmente en 2019 cuando llegué al Nacional y empecé forjar mis más grandes grupos de amigos. Ese camino me condujo a esta facultad, a esta carrera, y a este trabajo final. Me gustaría dar máxima gratitud a Patricia, Flora, Matías, Silvia, Gabriel y Verónica.

\vspace{\baselineskip}

\noindent Quiero agradecer a toda mi familia: especialmente a mis padres Gustavo y Guillermina, por su apoyo; y a mis abuelos Carlos, Ana, Lorenzo y Marta. Sé que los árboles genealógicos Ruiz Díaz y Herrera van a estar orgullosos de tener un nuevo universitario.

\vspace{\baselineskip}

\noindent Todas estas personas, y seguro muchas más, son las que me trajeron hasta acá: nunca estuve solo. Pido sinceras disculpas si me olvido de alguien. Quiero mencionar especialmente a Federico Brisnicoff, Tomás Molezzi, Santiago Manzanares, Carolina Castellucci, Franco Lotti, Sebastián Naddeo, Antonella Berzotti, Martina Ranieri, Facundo Álvarez Motta, Martina Álvarez Motta, Joaquín Fernández, Pilar Palmarino, Lucas Sandleris, Nicolás Schaievitch, Lisandro Acuña, Olivia López Taiana, Francisco Cirelli, Matteo Ginhson, Nicolás Pérez Stefan, Tomás Castro, Sofía Kutianski, Lucas Blake, Dante Culaciati, Nicolás Yañez, el Chino Cribioli, Franco Bongiovanni, Franco de Rosa, Facundo Decroix, Oriana Martínez y Franco Bottazzi.

\vspace{\baselineskip}

\noindent También quiero mencionar expresamente a Matías y Bruno, mi director y codirector respectivamente, ya que sin ellos, esta tesis no hubiera sido posible o su calidad hubiera sido menor.

\vspace{\baselineskip}

\noindent Por último, quiero agradecer a la comunidad de la cooperativa Eryx por estos diez meses maravillosos que pasé junto a ellos, previo a mi emigración a México, y en los cuales pude desarrollar el proyecto que hoy presento como Tesis de Licenciatura.
